% section: 極限

\section{極限}\label{節：極限}
  \noindent\textbf{【本編で使う場面】}\par
  数列や関数が近づく値を表す記法と計算を確認する.\sectionref*{節：ニュートン法型},発展演習の極限問題,付録で使う.基本・応用パターンの一般項を求める段階では後回しにできる.

  \subsection{極限}\label{節：極限：極限}
    ある関数$f(x)$について\,$x$を$a$に限りなく近づける操作を数学的には
    \[
      \lim_{x\to a} f(x)
    \]
    と表す.
    これを用いることで,ある関数$f(x)$の$x=a$の近傍での値がどうなっているかを知ることができる.
    例えば,\,$f(x)=x^2-4$のとき
    \[
      \lim_{x\to 2} (x^2-4)
    \]
    は$0$となる.
    一見シンプルな代入計算に見えるかもしれないが,これが今後重要な役割を果たす.

  \subsection{数列の極限}\label{節：数列の極限}
    数列$\{a_n\}$の番号$n$を大きくしていくとき,\,$a_n$がある値$L$に限りなく近づくなら,\,$a_n$は$L$に収束するといい,
    \[
      \lim_{n\to\infty}a_n=L
    \]
    と書く.
    
    例えば,\,$n>0$のとき
    \[
      a_n=\frac{1}{n}
    \]
    とすると,\,$n$を大きくするほど$a_n$は$0$に近づくので,
    \[
      \lim_{n\to\infty}\frac{1}{n}=0
    \]
    である.

  \subsection{関数の極限}\label{節：関数の極限}
    $x$を$a$に近づけたとき,\,$f(x)$が$L$に近づくなら,
    \[
      \lim_{x\to a}f(x)=L
    \]
    と書く.
    このとき,\,$x=a$での関数の値そのものではなく,\,$a$の近くでの値の変化に注目する.

    極限の計算では,既に求めた極限を組み合わせたり,不等式で上下から挟んだりする.
    例えば,
    \[
      a_n\le b_n\le c_n
    \]
    が成り立っており,
    \[
      \lim_{n\to\infty}a_n=\lim_{n\to\infty}c_n=L
    \]
    であるとき,\,$b_n$も$L$に収束する.
    これをはさみうちの原理という.

  \subsection{片側極限とその記法}\label{節：片側極限}
    $x$を$a$に近づけるとき,近づく側を限定した極限を\textbf{片側極限}という.
    $x<a$の範囲で$a$に近づける場合を\textbf{左側極限},\,$x>a$の範囲で$a$に近づける場合を\textbf{右側極限}という.
    数直線で考えると,それぞれ$a$の左側・右側から近づけることに対応する.

    左側極限には,次のような複数の記法がある.
    \[
      \boxed{
        \begin{aligned}
          \lim_{x\to a-}f(x)
          &=\lim_{x\to a^-}f(x)\\
          &=\lim_{\substack{x\to a\\x<a}}f(x)\\
          &=\lim_{x\uparrow a}f(x)
        \end{aligned}
      }
    \]
    これらは同じ左側極限を表す記法である.
    $a^-$の$-$は,\,$x<a$の側から近づけることを示す.
    また,\,$x\uparrow a$も$a$より小さい側から$a$に近づけることを表す.

    同様に,右側極限には次の記法がある.
    \[
      \boxed{
        \begin{aligned}
          \lim_{x\to a+}f(x)
          &=\lim_{x\to a^+}f(x)\\
          &=\lim_{\substack{x\to a\\x>a}}f(x)\\
          &=\lim_{x\downarrow a}f(x)
        \end{aligned}
      }
    \]
    $a^+$の$+$は,\,$x>a$の側から近づけることを示す.
    また,\,$x\downarrow a$も$a$より大きい側から$a$に近づけることを表す.
    特に$a=0$のときは,\,$x\to0^+$を$x\to+0$,\,$x\to0^-$を$x\to-0$と書くこともある.
    $+0$や$-0$は近づく側を表す記法であり,異なる二つの数を表しているわけではない.

    左右の両側から近づけることができる場合,通常の極限$\lim_{x\to a}f(x)$が有限の値$L$になるための必要十分条件は,左右の片側極限がともに存在し,同じ値$L$になることである.
    つまり,
    \[
      \lim_{x\to a^-}f(x)
      =\lim_{x\to a^+}f(x)=L
    \]
    である.
    例えば,
    \[
      f(x)=
      \begin{cases}
        -1 & (x<0),\\
        1 & (x\ge0)
      \end{cases}
    \]
    の場合,左側極限は$-1$,右側極限は$1$となる.
    左右の値が異なるので,\,$\lim_{x\to0}f(x)$は存在しない.
    また,片側極限でも$x=a$での関数の値そのものは関係しない.

  \subsection{\texorpdfstring{ネイピア数$e$の定義}{ネイピア数eの定義}}\label{節：ネイピア数eの定義}
    ネイピア数$e$は,ベルヌーイの定義によって,
    \[
      e=\lim_{n\to\infty}\left(1+\frac{1}{n}\right)^n
    \]
    と定める.
    以下では,この定義と数列・関数の極限を用いて,オイラーの公式の導出に必要な関係を確認する.

  \subsection{三角関数の基本極限}\label{節：三角関数の基本極限}
    角度は弧度法で表す.
    $0<x<\dfrac{\pi}{2}$のとき,半径$1$の円の扇形と,その扇形に内接・外接する三角形の面積を比べると,
    \begin{center}
      \BookMathLayout{\begin{tikzpicture}[
        x=2.4cm,
        y=2.4cm,
        line cap=round,
        line join=round,
        >=stealth
      ]
        % x=40度を例に,扇形と内接・外接する三角形を描く.
        \pgfmathsetmacro{\figangle}{40}
        \coordinate (O) at (0,0);
        \coordinate (A) at (1,0);
        \coordinate (P) at ({cos(\figangle)},{sin(\figangle)});
        \coordinate (Q) at (1,{tan(\figangle)});

        % 第一象限だけを描き,各図形は境界線で示す.
        \draw[->,thin] (O)--(1.38,0) node[below] {$x$};
        \draw[->,thin] (O)--(0,1.28) node[left] {$y$};
        \draw[thin] (A) arc[start angle=0,end angle=90,radius=1];
        \draw[thick] (O)--(Q)--(A)--cycle;
        \draw[thick] (O)--(P)--(A);
        \draw[thick] (A) arc[start angle=0,end angle=\figangle,radius=1];
        \draw[->,semithick] (0.24,0)
          arc[start angle=0,end angle=\figangle,radius=0.24];

        \fill (O) circle (0.018);
        \fill (A) circle (0.018);
        \fill (P) circle (0.018);
        \fill (Q) circle (0.018);
        \node[below left] at (O) {$O$};
        \node[below] at (0.52,-0.03) {$1$};
        \node[below right] at (A) {$A$};
        \node[above left] at (P) {$P$};
        \node[above right] at (Q) {$Q$};
        \node at (0.55,0.20) {$\frac12\sin x$};
        \node at (0.27,0.07) {$x$};

        \node[anchor=east,align=right,text width=2.0cm] (sector-label)
          at (-0.10,0.88) {扇形\\[-1pt]面積 $\frac12x$};
        \draw[->,thin] (sector-label.east)--(0.94,0.342);
        \node[anchor=west,align=left,text width=2.1cm] (outer-label)
          at (1.33,0.62) {外接三角形\\[-1pt]面積 $\frac12\tan x$};
        \draw[->,thin] (outer-label.west)--(0.91,0.55);
      \end{tikzpicture}}{\resizebox{\linewidth}{!}{\begin{tikzpicture}[
        x=2.4cm,
        y=2.4cm,
        line cap=round,
        line join=round,
        >=stealth
      ]
        % x=40度を例に,扇形と内接・外接する三角形を描く.
        \pgfmathsetmacro{\figangle}{40}
        \coordinate (O) at (0,0);
        \coordinate (A) at (1,0);
        \coordinate (P) at ({cos(\figangle)},{sin(\figangle)});
        \coordinate (Q) at (1,{tan(\figangle)});

        % 第一象限だけを描き,各図形は境界線で示す.
        \draw[->,thin] (O)--(1.38,0) node[below] {$x$};
        \draw[->,thin] (O)--(0,1.28) node[left] {$y$};
        \draw[thin] (A) arc[start angle=0,end angle=90,radius=1];
        \draw[thick] (O)--(Q)--(A)--cycle;
        \draw[thick] (O)--(P)--(A);
        \draw[thick] (A) arc[start angle=0,end angle=\figangle,radius=1];
        \draw[->,semithick] (0.24,0)
          arc[start angle=0,end angle=\figangle,radius=0.24];

        \fill (O) circle (0.018);
        \fill (A) circle (0.018);
        \fill (P) circle (0.018);
        \fill (Q) circle (0.018);
        \node[below left] at (O) {$O$};
        \node[below] at (0.52,-0.03) {$1$};
        \node[below right] at (A) {$A$};
        \node[above left] at (P) {$P$};
        \node[above right] at (Q) {$Q$};
        \node at (0.55,0.20) {$\frac12\sin x$};
        \node at (0.27,0.07) {$x$};

        \node[anchor=east,align=right,text width=2.0cm] (sector-label)
          at (-0.10,0.88) {扇形\\[-1pt]面積 $\frac12x$};
        \draw[->,thin] (sector-label.east)--(0.94,0.342);
        \node[anchor=west,align=left,text width=2.1cm] (outer-label)
          at (1.33,0.62) {外接三角形\\[-1pt]面積 $\frac12\tan x$};
        \draw[->,thin] (outer-label.west)--(0.91,0.55);
      \end{tikzpicture}}}
      \BookFigureCaption{扇形と内接・外接三角形の面積比較}
    \end{center}
    \[
      \frac{1}{2}\sin x
      <
      \frac{1}{2}x
      <
      \frac{1}{2}\tan x
    \]
    となる.
    したがって,
    \[
      \cos x<\frac{\sin x}{x}<1
    \]
    である.
    $x\to+0$のとき$\cos x\to1$なので,はさみうちの原理より
    \[
      \lim_{x\to+0}\frac{\sin x}{x}=1
    \]
    を得る.
    また,\,$\dfrac{\sin x}{x}$は偶関数であるから,左側から近づけた場合にも同じ極限となる.
    よって,
    \[
      \lim_{x\to0}\frac{\sin x}{x}=1
    \]
    である.

    ここで示した面積比較による証明は,初等的でよく用いられる説明である.
    ただし,扇形の面積公式や図形の大小関係を直観的に用いており,それらの根拠を厳密に示していないため,厳密な証明ではなく直観的な説明として位置づける.

    さらに,\,$\tan x=\dfrac{\sin x}{\cos x}$および$\lim_{x\to0}\cos x=1$より,
    \[
      \lim_{x\to0}\frac{\tan x}{x}
      =
      \lim_{x\to0}
      \left(
        \frac{\sin x}{x}\cdot\frac{1}{\cos x}
      \right)
      =1
    \]
    となる.

  \subsection{自然対数と指数関数の基本極限}\label{節：自然対数と指数関数の基本極限}
    以下では,\,$\log x$を自然対数$\log_e x$の意味で用いる.
    \[
      \lim_{x\to0}\frac{\log(1+x)}{x}
    \]
    について考える.
    $x$が$0$に十分近く,\,$x\ne0$のとき,対数の性質から
    \[
      \frac{\log(1+x)}{x}
      =\log\left((1+x)^\frac{1}{x}\right)
    \]
    である.右辺の極限を求めるには,
    \[
      \lim_{x\to0}(1+x)^\frac{1}{x}=e
    \]
    を使う.これは,先に指数対数の節で定めた$e$の極限表示から確認できる.
    $x\to 0^+$では,\,$\displaystyle t=\frac{1}{x}\to\infty$と置くと
    \[
      (1+x)^\frac{1}{x}=\left(1+\frac1t\right)^t\longrightarrow e.
    \]
    $x\to 0^-$では,\,$\displaystyle t=-\frac{1}{x}\to\infty$と置くと
    \begin{align*}
      &(1+x)^\frac{1}{x}
      =\left(1-\frac1t\right)^{-t}\\
      &=\left(1+\frac1{t-1}\right)^{t-1}
       \left(1+\frac1{t-1}\right)
      \longrightarrow e.
    \end{align*}
    自然対数は正の数上で連続であり,\,$\log e=1$なので,
\BookMathLayout{\[
      \lim_{x\to0}\frac{\log(1+x)}{x}
      =\lim_{x\to0}\log\left\{(1+x)^\frac{1}{x}\right\}
      =\log e=1
    \]}{
\[
\begin{aligned}
\lim_{x\to0}\frac{\log(1+x)}x
&=\lim_{x\to0}\log\left\{(1+x)^{1/x}\right\}\\
&=\log e=1
\end{aligned}
\]
}
    となる.
    この極限を用いて,\,$\lim_{n\to\infty}\left(1+\frac{a}{n}\right)^n$を求める.
    $a=0$の場合は,
    \[
      \lim_{n\to\infty}\left(1+\frac{0}{n}\right)^n=1=e^0
    \]
    である.以下,\,$a\ne0$とする.
    $n$が十分大きければ,\,$1+\dfrac{a}{n}>0$なので,
    \[
      \lim_{n\to\infty}\left(1+\frac{a}{n}\right)^n
      =
      \left\{
        \lim_{n\to\infty}
        \left(1+\frac{a}{n}\right)^{\frac{n}{a}}
      \right\}^{a}
      =e^a
    \]
    となる.以上より,\,$a=0$を含む任意の実数$a$について,
    \[
      \lim_{n\to\infty}\left(1+\frac{a}{n}\right)^n=e^a
    \]
    が成り立つ.
